\section{硬件与系统共设计：从 OCP 到 DGX}

前一节说明故障与恢复受系统形态影响，本节进一步追问计算节点本身为何变化。随着单节点集成更多 GPU、高速互联和专用网络，传统“小而可替换”的服务器哲学与“高密度、低通信延迟”的 AI 需求发生冲突，软件调度必须重新理解物理拓扑。

\subsection{架构回摆：从分散微服务到高密度计算节点}
这一回摆不是简单复古，而是 workload 对通信提出的新要求。课程把 Cray、OCP 与现代 AI box 放在同一历史轴上，帮助读者比较节点密度、可维护性和互联效率；读图重点是设计目标怎样变化，而不是把不同年代设备做性能排名。
讲者回顾了 Cray-2、OCP 服务器与现代 AI 机箱的演化。传统云强调模块化小节点，利于替换和弹性；AI 训练/推理要求高带宽互联与低通信延迟，又推动架构向高密度节点回摆。

\begin{importantbox}{AI 负载下 ``模块化'' 与 ``集成化'' 的再平衡}
模块化降低运维复杂度；集成化提升互联效率。AI 系统常需要在集成节点内部追求极致带宽，再在节点间做分层调度。
\end{importantbox}

\subsection{节点形态对软件栈的影响}
当单机从 ``若干独立 CPU'' 变为 ``8 GPU + 2 CPU + 高速互联''，软件抽象也会变化：
\begin{itemize}
    \item 资源单位不再只是 ``CPU 核 + 内存''，而是 ``拓扑感知的加速器组''；
    \item 调度器需要理解 NVLink / PCIe / NIC 布局；
    \item 容器编排需要加入通信亲和性约束。
\end{itemize}

\begin{knowledgebox}{为什么硬件细节重新进入应用层视野}
过去很多应用团队不需要关心机架与交换机拓扑；在大模型任务里，物理拓扑直接影响训练时延、吞吐和成本。系统抽象不能完全屏蔽底层结构。
\end{knowledgebox}

\begin{warningbox}{过度抽象的代价}
如果平台把所有异构资源都抽成 ``统一 VM''，短期易用性更好，长期会掩盖通信瓶颈与拓扑冲突，导致成本优化空间被锁死。
\end{warningbox}

\subsection{本章小结}
AI 基础设施正在经历 ``硬件-软件共设计'' 回归。高密度节点、拓扑感知调度和分层抽象将成为主流设计方向。

